\subsection{A limitation of the auxiliary level-twenty-seven Kummer condition}
\label{subsec:auxiliary-level27-kummer}

One possible auxiliary test uses a rational three-torsion family
and the rank-zero curve \(X_0(27)\). The following explicit
calculation shows that its integral three-adic Kummer condition
fails to distinguish the boundary at infinity from either rational
modular cusp class, on a positive-density set of good primes.
The conclusion concerns this auxiliary condition; it supplies no
global Selmer counterexample and does not decide the full compact
Fermat point-map fibre discussed above.

Write
\[
 C=X_0(27):\quad y^2+y=x^3-7,\qquad
 h=-3+\frac{y+5}{x},\qquad
 \lambda=\left(\frac{h}{h+9}\right)^3.
\]
The model and the map \(h:C\to X_0(9)\) are given in
\cite[proof of Proposition 24, p.~18 of the 2020 author version]{BourdonGillRouseWatson2024}.
Their displayed \(j\)-map, followed by direct substitution, gives
\[
 j=\frac{(h+9)^3(h^3+243h^2+2187h+6561)^3}
          {h^9(h^2+9h+27)}
   =\frac{27(9-8\lambda)^3}{\lambda^3(1-\lambda)}.
\]
Thus this is the relevant level-\(27\) cover for the family
\[
 \mathcal E_\lambda:\quad
 Y^2+XY+\frac{\lambda}{27}Y=X^3.
\]
This rational \(3\)-torsion family is recalled in
\cite[\S2.3, (2.3.3) and Lemma 2.3.5, p.~7]{FreiHonigsVoight2025}.
Its marked point \((0,0)\) fixes a vector modulo three, while the
geometric determinant is one. Thus the residual geometric image
is contained in the unipotent group of order three. The higher
congruence kernel is pro-\(3\), so the entire geometric Tate image
is pro-\(3\). This makes the family available on the geometric
pro-\(3\) quotient, but does not make its auxiliary Kummer condition
separate the boundary points.
Let \(O\) be the origin of \(C\), and put \(T=(3,4)\).
These are rational modular cusps, with \(\lambda(O)=1\) and
\(\lambda(T)=0\). The third rational point \((3,-5)\) is
noncuspidal: it has \(\lambda=-1/8\) and \(j=-12288000\).
The distinction between the two rational modular cusps and the
rational noncuspidal self-point is recorded by
\citet[\S6.2, p.~17 of the 17 March 2008 author version]{WuthrichSelfPoints2009}.
Put \(D_\infty=\lambda^{-1}(\infty)\).
For \(q\ne3\), write
\[
 \kappa_q:C(\Q_q)\longrightarrow H^1(\Q_q,T_3C)
\]
for the integral \(3\)-adic Kummer map based at \(O\).

\begin{lemma}[Failure of auxiliary Kummer separation]
\label{lem:auxiliary-level27-kummer-failure}
For each \(A\in\{O,T\}\), there is a positive-density set of rational
primes \(q\ne3\) for which some \(R_q\in D_\infty(\Q_q)\) satisfies
\[
 \kappa_q(R_q)=\kappa_q(A).
\]
The equality also holds on a \(q\)-adic neighborhood of \(R_q\) in
\(C(\Q_q)\). Explicit witnessing primes are \(q=17\) for \(A=O\)
and \(q=71\) for \(A=T\).
\end{lemma}

\begin{proof}
The discriminant of \(C\) is \(-3^9\), so \(C\) has good reduction at
every \(q\ne3\). Its reduction kernel is a pro-\(q\) group on which
multiplication by \(3\) is an automorphism. Consequently
\[
 C(\Q_q)^{\wedge3}\simeq C(\F_q)^{\wedge3},
\]
and equality of the \(3\)-primary reduction classes is equivalent to
equality under \(\kappa_q\).

Substitution of \(y=x(h+3)-5\) in the equation of \(C\) gives
\[
 F(x,h)=x^3-(h+3)^2x^2+(9h+27)x-27=0.
\]
The fibre \(D_\infty\) is \(h=-9\); hence its points have
\begin{equation}\label{eq:auxiliary-level27-fibre}
 f(x)=x^3-36x^2-54x-27=0,\qquad y=-6x-5.
\end{equation}
Here \(\operatorname{disc}(f)=-3^{13}\).
Modulo \(17\), equation~\eqref{eq:auxiliary-level27-fibre} has the point
\[
 \overline R=(12,8).
\]
It is nonzero \(2\)-torsion, since \(2\cdot8+1=0\) in \(\F_{17}\).
Also \(f'(12)=7\pmod {17}\). Hensel's lemma therefore supplies a
point \(R_{17}\in D_\infty(\Q_{17})\) reducing to \(\overline R\).
Its \(3\)-primary reduction class is zero, proving the claim for
\(A=O\) at \(17\).

Modulo \(71\), the point \(\overline R=(69,7)\) belongs to~\eqref{eq:auxiliary-level27-fibre}, and
\(f'(69)=31\pmod {71}\). With \(\overline D=\overline R-\overline T\),
the group law gives
\[
 \overline D=(36,69),\quad
 2\overline D=(39,55),\quad
 4\overline D=(38,35),\quad
 8\overline D=O.
\]
Thus the Hensel lift \(R_{71}\in D_\infty(\Q_{71})\) has the same
\(3\)-adic Kummer class as \(T\). Constancy on the corresponding
residue neighborhoods follows from the reduction comparison.

We give the finite-extension argument for positive density explicitly.
Fix one of these pairs \((q_0,A)\), and choose the algebraic point
\(R\) of~\eqref{eq:auxiliary-level27-fibre} with the specified embedding into \(\Q_{q_0}\).
Set \(K=\Q(R)\). Since \(f\) is monic with discriminant a power of
\(3\), \(K\) is unramified away from \(3\).
The point \(R-A\) is divisible by \(9\) in \(C(\Q_{q_0})\):
its reduction has order dividing \(8\), and multiplication by \(9\)
is an automorphism on the reduction kernel. Choose an algebraic
point \(B\), embedded in \(\Q_{q_0}\), such that
\begin{equation}\label{eq:auxiliary-level27-division}
 [9]B=R-A.
\end{equation}
Let \(L/\Q\) be a finite Galois extension containing
\(K\), the coordinates of \(B\), \(C[9]\), and \(\mu_{27}\).
It may be chosen unramified away from \(3\). Indeed, \(C\) is an
abelian scheme away from \(3\), every \(K\)-point extends there by
properness, and multiplication by \(9\) is finite \'etale there.
Adjoining the finite \'etale fibre above \(R-A\), taking its Galois
closure, and adjoining \(\mu_{27}\) introduce no ramification
away from \(3\).

Choose a prime of \(L\) above \(q_0\), and let \(\sigma\) be its
arithmetic Frobenius. Since the chosen \(R\) and \(B\) are
\(\Q_{q_0}\)-rational, \(\sigma\) fixes both. Its action on
\(\mu_{27}\) is the \(17\)-th power: both \(17\) and \(71\) are
\(17\pmod {27}\).
Chebotarev's theorem
\cite[Chapter I, \S2.2, pp.~I-7--I-8]{Serre1989}
gives a positive-density set of primes whose Frobenius class in
\(\operatorname{Gal}(L/\Q)\) is the class of \(\sigma\).
For each such \(q\), choose a prime above \(q\) with Frobenius
exactly \(\sigma\). The extension is unramified, so its decomposition
group is generated by \(\sigma\). Therefore \(R\) and \(B\) give
points in \(C(\Q_q)\), and~\eqref{eq:auxiliary-level27-division} holds there. Moreover
\begin{equation}\label{eq:auxiliary-level27-congruence}
 q\equiv17\pmod {27}.
\end{equation}

For \(q\equiv2\pmod3\), cubing is a bijection on \(\F_q\).
For each \(y\in\F_q\), the equation \(x^3=y^2+y+7\) has exactly one
solution, whence \(\#C(\F_q)=q+1\).
By~\eqref{eq:auxiliary-level27-congruence}, the \(3\)-primary subgroup has order exactly \(9\).
Equation~\eqref{eq:auxiliary-level27-division} consequently makes the \(3\)-primary reduction class
of \(R-A\) zero. The good-reduction comparison proves the required
Kummer equality for every prime in this positive-density set.
\end{proof}

\begin{remark}[Local Fermat points and the remaining scope]
\label{rem:auxiliary-level27-local-fermat}
These collisions persist on actual local points of the level cover of
\(U_a\), for \(a=5\) or \(7\).
At any of the preceding primes \(q\ne a\), take \(w=q^{-m}\) with
\(m\ge1\), and set
\[
 \lambda=w^{3a},\qquad
 h=\frac{9w^a}{1-w^a}.
\]
As \(m\) increases, \(h\) approaches \(-9\).
Since \(f'\) is a unit at the chosen residue root, the implicit
function theorem gives a point of \(C(\Q_q)\) near \(R_q\) with
this value of \(h\). Its Kummer class remains that of \(A\).
At the same time,
\[
 x_F=w^3,\qquad
 y_F=-w^3(1-w^{-3a})^{1/a}
\]
are defined over \(\Q_q\) by Hensel's lemma and satisfy
\(x_F^a+y_F^a=1\), hence give the points
\([x_F:y_F:-1]\) of the all-plus model of \(U_a\).
They reduce to the coordinate cusp \((1:-1:0)\), where
\(\lambda=\infty\).
The modular point supplies the cyclic level-\(27\) subgroup for
\(\mathcal E_\lambda\); away from the exceptional \(j\)-values,
quadratic twisting does not change the field of definition of such
a subgroup.

Thus the auxiliary \(X_0(27)\) Kummer condition does not separate
this boundary from either rational cusp class, even after deleting
a finite set of primes. This is a calculation about that condition
only. It does not construct a global Selmer section, and it does
not identify the full compact pro-$3$ section of \(X_a\) at this
boundary with the prescribed rational cusp. The latter finite-field
fibre condition remains an additional, unresolved requirement.
\end{remark}

